% GENERATED FILE. Edit canonical lessons, then run npm run generate:books.
% canonical-source-hash: fnv1a64:32e7f4525a40acea
% canonical-lessons: GU-C214-padvu, GU-C214-dhodhmar, GU-C214-tipu, GU-C214-tapakvu, GU-C214-palalvu

\chapter{Rain on the way}
\label{ch:gu-rain-on-the-way}
\hlchaptermodality{214}

\begin{chapteropening}
\textbf{By the end of this chapter:} \emph{I can say that rain will fall and that I will get wet.}

Complete the last lesson of chapter 214: I can say that rain will fall and that I will get wet.
\end{chapteropening}

\section[\texorpdfstring{\gu{પડવું} (paḍvũ)}{પડવું (paḍvũ)}]{\gu{પડવું} (paḍvũ) — to fall, to happen}

\label{lesson:GU-C214-padvu}



\begin{quote}
Before the new one: say the Gujarati for cloudy, then the Gujarati for to pour down.
\end{quote}

\subsection*{You'll want to know: \gu{પડવું}}
\textbf{\gu{પડવું}} — \emph{paḍvũ} — \textquotedblleft{}to fall, to happen\textquotedblright{}.

\begin{quote}
\textbf{\gu{સાંજે} \gu{વરસાદ} \gu{પડશે}.} — \emph{sā̃je varsād paḍśe} — \textquotedblleft{}It will rain in the evening.\textquotedblright{}
\end{quote}

\begin{cousinweb}
\textbf{\gu{પડવું}} (\emph{paḍvũ}) is a tadbhava. It comes from Sanskrit \emph{patati}, \textquotedblleft{}flies; falls\textquotedblright{}, and CDIAL glosses the Gujarati verb \textquotedblleft{}to fall, happen\textquotedblright{} (CDIAL 7722).
\end{cousinweb}

\begin{grammarlens}[title={What you've built}]
One more word for everyday life at A2.
\end{grammarlens}

\subsection*{Your turn}
\begin{itemize}
\raggedright
  \item \emph{Say it:} \emph{paḍvũ}
  \item \emph{Say it:} \emph{paḍvũ}, once more
  \item \emph{Say it:} \emph{sā̃je varsād paḍśe} — It will rain in the evening.
  \item \emph{Recall:} say the Gujarati for cloudy, then the Gujarati for to pour down, then say \emph{paḍvũ} again
\end{itemize}

\begin{tcolorbox}[breakable,colback=teal!4,colframe=teal!35!black,title={Before you move on}]
What does \gu{પડવું} mean? (\textquotedblleft{}To fall, to happen\textquotedblright{}.) Say it once more.
\end{tcolorbox}
\section[\texorpdfstring{\gu{ધોધમાર} (dhodhmār)}{ધોધમાર (dhodhmār)}]{\gu{ધોધમાર} (dhodhmār) — pouring, torrential}

\label{lesson:GU-C214-dhodhmar}



\begin{quote}
Before the new one: say the Gujarati for to pour down, then the Gujarati for to fall.
\end{quote}

\subsection*{You'll want to know: \gu{ધોધમાર}}
\textbf{\gu{ધોધમાર}} — \emph{dhodhmār} — \textquotedblleft{}pouring, torrential\textquotedblright{}.

\begin{quote}
\textbf{\gu{કાલે} \gu{ધોધમાર} \gu{વરસાદ} \gu{પડશે}.} — \emph{kāle dhodhmār varsād paḍśe} — \textquotedblleft{}There will be very heavy rain tomorrow.\textquotedblright{}
\end{quote}

It does not change for gender.

\begin{cousinweb}
\textbf{\gu{ધોધમાર}} (\emph{dhodhmār}) is built on \emph{dhodh}, \textquotedblleft{}waterfall\textquotedblright{}, so the rain comes down like a waterfall. The second part, \emph{-mār}, is often linked to \emph{mārvũ}, \textquotedblleft{}to strike\textquotedblright{}, but CDIAL lists no form of the compound, so its make-up is not settled.
\end{cousinweb}

\begin{grammarlens}[title={What you've built}]
One more word for everyday life at A2.
\end{grammarlens}

\subsection*{Your turn}
\begin{itemize}
\raggedright
  \item \emph{Say it:} \emph{dhodhmār}
  \item \emph{Say it:} \emph{dhodhmār}, once more
  \item \emph{Say it:} \emph{kāle dhodhmār varsād paḍśe} — There will be very heavy rain tomorrow.
  \item \emph{Recall:} say the Gujarati for to pour down, then the Gujarati for to fall, then say \emph{dhodhmār} again
\end{itemize}

\begin{tcolorbox}[breakable,colback=teal!4,colframe=teal!35!black,title={Before you move on}]
What does \gu{ધોધમાર} mean? (\textquotedblleft{}Pouring, torrential\textquotedblright{}.) Say it once more.
\end{tcolorbox}
\section[\texorpdfstring{\gu{ટીપું} (ṭīpũ)}{ટીપું (ṭīpũ)}]{\gu{ટીપું} (ṭīpũ) — a drop}

\label{lesson:GU-C214-tipu}



\begin{quote}
Before the new one: say the Gujarati for to fall, then the Gujarati for pouring.
\end{quote}

\subsection*{You'll want to know: \gu{ટીપું}}
\textbf{\gu{ટીપું}} — \emph{ṭīpũ} — \textquotedblleft{}a drop\textquotedblright{}.

\begin{quote}
\textbf{\gu{પાણીનું} \gu{એક} \gu{ટીપું} \gu{પડશે}.} — \emph{pāṇīnũ ek ṭīpũ paḍśe} — \textquotedblleft{}A drop of water will fall.\textquotedblright{}
\end{quote}

\begin{cousinweb}
\textbf{\gu{ટીપું}} (\emph{ṭīpũ}) comes from \emph{*ṭipp}, \textquotedblleft{}drip, drop\textquotedblright{}, one of a set of imitative reconstructed forms that CDIAL gathers under \emph{*ṭapp}; it lists Gujarati \emph{ṭīpũ}, \textquotedblleft{}drop\textquotedblright{} (CDIAL 5444).
\end{cousinweb}

\begin{grammarlens}[title={What you've built}]
One more word for everyday life at A2.
\end{grammarlens}

\subsection*{Your turn}
\begin{itemize}
\raggedright
  \item \emph{Say it:} \emph{ṭīpũ}
  \item \emph{Say it:} \emph{ṭīpũ}, once more
  \item \emph{Say it:} \emph{pāṇīnũ ek ṭīpũ paḍśe} — A drop of water will fall.
  \item \emph{Recall:} say the Gujarati for to fall, then the Gujarati for pouring, then say \emph{ṭīpũ} again
\end{itemize}

\begin{tcolorbox}[breakable,colback=teal!4,colframe=teal!35!black,title={Before you move on}]
What does \gu{ટીપું} mean? (\textquotedblleft{}A drop\textquotedblright{}.) Say it once more.
\end{tcolorbox}
\section[\texorpdfstring{\gu{ટપકવું} (ṭapakvũ)}{ટપકવું (ṭapakvũ)}]{\gu{ટપકવું} (ṭapakvũ) — to drip, to leak}

\label{lesson:GU-C214-tapakvu}



\begin{quote}
Before the new one: say the Gujarati for pouring, then the Gujarati for a drop.
\end{quote}

\subsection*{You'll want to know: \gu{ટપકવું}}
\textbf{\gu{ટપકવું}} — \emph{ṭapakvũ} — \textquotedblleft{}to drip, to leak\textquotedblright{}.

\begin{quote}
\textbf{\gu{વરસાદમાં} \gu{છત} \gu{ટપકશે}.} — \emph{varsādmā̃ chat ṭapakśe} — \textquotedblleft{}The roof will leak in the rain.\textquotedblright{}
\end{quote}

\begin{cousinweb}
\textbf{\gu{ટપકવું}} (\emph{ṭapakvũ}) is built on \emph{*ṭapp}, \textquotedblleft{}drip, drop\textquotedblright{}, an imitative reconstructed form. CDIAL lists the Gujarati noun \emph{ṭapkũ}, \textquotedblleft{}drop\textquotedblright{}, there, but not the verb itself (CDIAL 5444).
\end{cousinweb}

\begin{grammarlens}[title={What you've built}]
One more word for everyday life at A2.
\end{grammarlens}

\subsection*{Your turn}
\begin{itemize}
\raggedright
  \item \emph{Say it:} \emph{ṭapakvũ}
  \item \emph{Say it:} \emph{ṭapakvũ}, once more
  \item \emph{Say it:} \emph{varsādmā̃ chat ṭapakśe} — The roof will leak in the rain.
  \item \emph{Recall:} say the Gujarati for pouring, then the Gujarati for a drop, then say \emph{ṭapakvũ} again
\end{itemize}

\begin{tcolorbox}[breakable,colback=teal!4,colframe=teal!35!black,title={Before you move on}]
What does \gu{ટપકવું} mean? (\textquotedblleft{}To drip, to leak\textquotedblright{}.) Say it once more.
\end{tcolorbox}
\section[\texorpdfstring{\gu{પલળવું} (palaḷvũ)}{પલળવું (palaḷvũ)}]{\gu{પલળવું} (palaḷvũ) — to get soaked, to get wet}

\label{lesson:GU-C214-palalvu}



\begin{quote}
Before the new one: say the Gujarati for a drop, then the Gujarati for to drip.
\end{quote}

\subsection*{You'll want to know: \gu{પલળવું}}
\textbf{\gu{પલળવું}} — \emph{palaḷvũ} — \textquotedblleft{}to get soaked, to get wet\textquotedblright{}.

\begin{quote}
\textbf{\gu{હું} \gu{વરસાદમાં} \gu{પલળીશ}.} — \emph{hũ varsādmā̃ palaḷīś} — \textquotedblleft{}I will get soaked in the rain.\textquotedblright{}
\end{quote}

\begin{cousinweb}
\textbf{\gu{પલળવું}} (\emph{palaḷvũ}) is probably a tadbhava. CDIAL lists its transitive partner \emph{palāḷvũ}, \textquotedblleft{}to make wet, soak\textquotedblright{}, under \emph{*pralālayate}, \textquotedblleft{}slobbers\textquotedblright{}, a reconstructed form; \emph{palaḷvũ} itself is not listed (CDIAL 8759).
\end{cousinweb}

\begin{grammarlens}[title={What you've built}]
One more word for everyday life at A2.
\end{grammarlens}

\subsection*{Your turn}
\begin{itemize}
\raggedright
  \item \emph{Say it:} \emph{palaḷvũ}
  \item \emph{Say it:} \emph{palaḷvũ}, once more
  \item \emph{Say it:} \emph{hũ varsādmā̃ palaḷīś} — I will get soaked in the rain.
  \item \emph{Recall:} say the Gujarati for a drop, then the Gujarati for to drip, then say \emph{palaḷvũ} again
\end{itemize}

\begin{tcolorbox}[breakable,colback=teal!4,colframe=teal!35!black,title={Before you move on}]
What does \gu{પલળવું} mean? (\textquotedblleft{}To get soaked, to get wet\textquotedblright{}.) Say it once more.
\end{tcolorbox}
